\documentclass[11pt]{article}
\usepackage[T1]{fontenc}
\usepackage[utf8]{inputenc}
\usepackage{lmodern,amsmath,amssymb,amsthm,mathtools}
\usepackage[margin=28mm]{geometry}
\usepackage{microtype,xcolor}
\usepackage[authoryear,round]{natbib}
\usepackage{xurl}
\usepackage[colorlinks=true,allcolors=blue!50!black]{hyperref}
\newtheorem{theorem}{Theorem}[section]
\newtheorem{proposition}[theorem]{Proposition}
\newtheorem{lemma}[theorem]{Lemma}
\newtheorem{corollary}[theorem]{Corollary}
\theoremstyle{remark}\newtheorem{remark}[theorem]{Remark}
\newcommand{\Spec}{\operatorname{Spec}}
\newcommand{\OO}{\mathcal O}
\newcommand{\et}{\mathrm{\acute et}}
\newcommand{\Ext}{\mathcal E}
\newcommand{\BT}{\mathrm{BT}}
\newcommand{\workbenchurl}{https://github.com/KokunoYumeto/schemes-and-moduli-workbench}
\newcommand{\recordurl}{https://github.com/KokunoYumeto/mathematics-workbenches}
\title{A reduced counterexample to Grothendieck's extension assertion,\ after Ferrand--Raynaud and Olberding}
\author{}
\date{23 November 2026}
\begin{document}
\maketitle
\begin{abstract}
Grothendieck's notes on good reduction propose extending an abelian scheme together with a prescribed Barsotti--Tate group across a closed set of residue characteristic $p$. We construct a one-dimensional reduced geometrically unibranch noetherian local base on which this extension fails, even for an ordinary elliptic curve that itself has good reduction. Its normalization is not finite. The construction uses the derivation method for analytically ramified rings of Ferrand--Raynaud and Olberding, together with Serre--Tate theory. We also explain why a nilpotent thickening in characteristic $p$ preserves the extension property, while a simple mixed-characteristic thickening can destroy it. Thus reducedness alone neither detects the property nor repairs the proposed assertion. The mathematics and text were produced and checked by GPT-6 Astra Pro, including a separate statement-only re-proof by the same model; no human has checked them.
\end{abstract}

\section{Grothendieck's extension problem}

The question is whether an abelian scheme on an open subscheme is determined and extended by an integral model of its $p$-divisible group. Let $S$ be locally noetherian, let $U\subset S$ be open, and assume that every point of $T=S\setminus U$ has residue characteristic $p$. Define $\Ext(S,U)$ to have objects
\[
 (A_U,M,\varphi),\qquad
 \varphi:M|_U\xrightarrow{\sim}A_U[p^\infty],
\]
where $A_U/U$ is an abelian scheme and $M/S$ is a $p$-divisible group. A morphism is a homomorphism of abelian schemes on $U$ and a homomorphism of $p$-divisible groups on $S$ compatible with the displayed identifications. The extension functor is
\[
 F_{S,U}:\operatorname{Ab}(S)\longrightarrow\Ext(S,U),
 \qquad A\longmapsto(A|_U,A[p^\infty],\mathrm{id}).
\]
To realize a triple means to give an abelian scheme over $S$ together with isomorphisms realizing all three prescribed data.

In \emph{Bonne r\'eduction des vari\'et\'es ab\'eliennes via bonne r\'eduction de Barsotti--Tate}, Grothendieck writes
\begin{quote}
``Alors $F$ est une \'equivalence des cat\'egories''.
\end{quote}
The proposed statement assumes a finite surjective map $S'\to S$ with $S'$ geometrically unibranch and the depth condition
\begin{equation}\label{eq:depth}
 \operatorname{prof}^{\et}_T(S)\geq2.
\end{equation}
The passage is in cote~100, archive sheet~2 (handwritten page~1, PDF image~3) \citep{GrothendieckNotes}. For the local one-dimensional schemes below, \eqref{eq:depth} follows from the connectedness of the punctured strict local scheme \citep[Expos\'e~XIV, Remark~1.6(b), p.~164]{SGA2}. Equivalently, locally constant \'etale coefficients have no local cohomology in degrees zero and one supported at the closed point. These are \'etale conditions; they do not impose finiteness of normalization.

\begin{theorem}[A reduced counterexample]\label{thm:reduced}
There is a one-dimensional geometrically unibranch noetherian local domain $R$ of mixed characteristic $(0,5)$ with the following properties. Its normalization is a discrete valuation ring and is not finite over $R$. Set $S=\Spec R$ and $U=D(5)$; then $U$ is the generic point, the finite-cover hypothesis above holds with $S'=S$, and \eqref{eq:depth} holds. Nevertheless $F_{S,U}$ is not essentially surjective.

One may take
\[
 M=\mu_{5^\infty,R}\oplus(\mathbf Q_5/\mathbf Z_5)_R
\]
and an ordinary elliptic generic fibre $A_U$ that already extends to an elliptic scheme over $R$. The obstruction is to realizing the specified $p$-divisible group and generic identification together.
\end{theorem}

The example separates existence of some good model from extension of a specified integral comparison. It is relevant to arguments that normalize a base and then try to descend an infinite compatible system of torsion identifications. Although the normalization here is a universal homeomorphism, it is not a finite morphism. Finite-presentation descent statements cannot by themselves descend the entire system.

The opposite phenomenon occurs in characteristic $p$.

\begin{theorem}[A nonreduced base with the extension property]\label{thm:positive}
Let $k=\overline{\mathbf F}_p$, let
\[
 S=\Spec\bigl(k[[t]][\epsilon]/(\epsilon^2)\bigr),
 \qquad U=D(t).
\]
Then $F_{S,U}$ is an equivalence in every relative dimension. The two hypotheses of the proposed statement hold, although the local ring at $S\setminus U$ is nonreduced.
\end{theorem}

\paragraph{Credit and the known boundary.}
The derivation construction of an analytically ramified local domain is established commutative algebra. Ferrand and Raynaud introduced the method \citep[Lemma~2.1, pp.~300--303]{FerrandRaynaud}; \citet[Theorems~4.1 and~4.4, preprint pp.~11--13]{OlberdingBad} gives the differential-subring classification and the corresponding completion as an idealization. Our ring lemma spells out the rank-one mixed-characteristic case needed to prescribe the derivative of an elliptic $j$-invariant. The elliptic argument uses the Serre--Tate equivalence and ordinary canonical lifts in \citet[Theorem~1.2.1, pp.~143--145; Theorem~2.1, pp.~148--152]{Katz}. The good-reduction and homomorphism-extension inputs for the positive example are de Jong's theorems \citep[Theorems~2--3, p.~261]{DeJongICM}.

Normal bases satisfy a substantially stronger positive statement: over a locally noetherian normal base, with $U$ dense and the missing locus in characteristic $p$, extension of the $p$-divisible group is equivalent to extension of the abelian scheme \citep[Corollary~6.12]{KrishnamoorthyPal}. Theorem~\ref{thm:reduced} does not concern normal bases and does not modify that published theorem. The ring construction itself is known; its use here makes the missing descent requirement in Grothendieck's proposed assertion explicit even for reduced bases. No priority assertion is made for this application.

\paragraph{The mechanism.}
Start with a canonical ordinary elliptic lift, whose $p$-divisible group is split. A derivation turns its $j$-invariant $J$ into $J+\epsilon$ on the completion of a reduced subring. On the characteristic-$5$ dual-number fibre, the prescribed split group is constant. Serre--Tate theory then forces the elliptic deformation to be constant, while its $j$-invariant varies. The contradiction appears already modulo $5$. The infinite construction is used only to put this infinitesimal discrepancy inside a noetherian domain.

\section{Nilpotent thickenings in characteristic \texorpdfstring{$p$}{p}}

This section proves the positive example and identifies exactly the hypothesis under which reduction preserves the extension problem.

\begin{proposition}\label{prop:nilpotent}
Suppose that $p$ is locally nilpotent on the locally noetherian scheme $S$. Let $S_0\hookrightarrow S$ be a nilpotent closed immersion inducing the same underlying topological space, and put $U_0=U\times_SS_0$. A triple in $\Ext(S,U)$ is realizable if and only if its reduction in $\Ext(S_0,U_0)$ is realizable. If $F_{S_0,U_0}$ is fully faithful, so is $F_{S,U}$. In particular an equivalence on $S_0$ gives an equivalence on $S$.
\end{proposition}
\begin{proof}
Restriction proves one implication. For the other, choose a realization $A_0/S_0$ of the reduced triple, with its prescribed identification $A_0[p^\infty]\simeq M_{S_0}$. On an affine open of $S$, the Serre--Tate equivalence applies because both $p$ and the defining ideal of $S_0$ are nilpotent \citep[Theorem~1.2.1]{Katz}. It gives an abelian scheme $A$ with this reduction and with $p$-divisible group $M$. On $U$, its reduction and $p$-divisible data coincide with those of $A_U$. Full faithfulness of the same equivalence gives the unique isomorphism $A|_U\simeq A_U$ inducing the prescribed reduction and group comparison. On overlaps, uniqueness gives compatible isomorphisms, so the local schemes glue over $S$.

For morphisms of triples, first use full faithfulness over $S_0$ to obtain a unique homomorphism of their reductions. The given homomorphism of $p$-divisible groups lifts it uniquely by Serre--Tate. Its restriction to $U$ agrees with the prescribed homomorphism by the same uniqueness. This proves full faithfulness.
\end{proof}

\begin{proof}[Proof of Theorem~\ref{thm:positive}]
Let $S_0=\Spec k[[t]]$ and $U_0=\Spec k((t))$. For any triple over this trait, de Jong's good-reduction criterion gives an abelian scheme extending its generic abelian variety because the associated $p$-divisible group has a model over $S_0$ \citep[Theorem~3, p.~261]{DeJongICM}. His homomorphism-extension theorem extends the given generic isomorphism of $p$-divisible groups and its inverse uniquely \citep[Theorem~2, p.~261]{DeJongICM}. The extensions are inverse by faithfulness, hence the specified triple is realized. A homomorphism between the generic fibres of two abelian schemes over a discrete valuation ring extends uniquely, by the N\'eron mapping property of an abelian scheme \citep[Proposition~4.13 and Remark~4.14, author manuscript p.~13]{LiuTong}; compatibility with the prescribed $p$-divisible homomorphism follows from the same faithfulness. Thus $F_{S_0,U_0}$ is an equivalence.

Proposition~\ref{prop:nilpotent} applies to the square-zero thickening $S_0\hookrightarrow S$ and proves the equivalence over $S$. The finite closed immersion $S_0\to S$ is surjective and has normal, hence geometrically unibranch, source. The \'etale sites of $S$ and $S_0$ agree under a nilpotent thickening \citep[Tag~04DY]{Stacks}. The strictly henselian trait $S_0$ has a nonempty connected punctured spectrum. For any constant finite coefficient group, restriction on $H^0$ is therefore an isomorphism, and $H^1(S_0,-)=0$. The local-cohomology sequence gives the vanishing in degrees zero and one in \eqref{eq:depth}. Finally, $\epsilon\ne0$ is nilpotent in the local ring at the closed point.
\end{proof}

\section{The mixed-characteristic first-order obstruction}

A small explicit example shows why the nilpotence of $p$ cannot be omitted from Proposition~\ref{prop:nilpotent}.

\begin{proposition}\label{prop:mixed}
Let $R_1=\mathbf Z_5[\epsilon]/(\epsilon^2)$, let $S_1=\Spec R_1$, and let $U_1=D(5)$. Consider
\[
 E/\mathbf Z_5:\quad y^2=x^3+x+1,
 \qquad A_{U_1}:\quad y^2=x^3+x+(1+\epsilon).
\]
Set $M=E[5^\infty]_{S_1}$. There is a unique isomorphism
$M|_{U_1}\simeq A_{U_1}[5^\infty]$ reducing to the identity modulo $\epsilon$. The resulting triple is not realizable over $S_1$, although the finite-cover and depth hypotheses in Section~1 hold.
\end{proposition}
\begin{proof}
The discriminant of $E$ is $-16(4+27)=-496$, a $5$-adic unit. Hence $E$ is an elliptic scheme. Over $U_1$, the prime $5$ is invertible, so all the $5^m$-torsion group schemes are finite \'etale. Invariance of finite \'etale schemes under a nilpotent thickening gives the unique identity-preserving isomorphism at each level. Uniqueness makes these isomorphisms compatible in $m$.

Suppose an abelian scheme $A/S_1$ realizes the triple. Its reduction modulo $\epsilon$ and $E$ have isomorphic generic fibres over $\mathbf Q_5$, so they are isomorphic over $\mathbf Z_5$ by uniqueness of good reduction. The reduction of the specified group comparison agrees with this isomorphism: both agree generically, and homomorphisms of $p$-divisible groups over a trait are determined by their generic fibres. For each $r\geq1$, Serre--Tate over $(\mathbf Z_5/5^r)[\epsilon]/(\epsilon^2)$ now identifies $A$ with the constant deformation of $E$, because their reduction and prescribed $p$-divisible group agree. Thus
\[
 j(A)-j(E)\in\bigcap_{r\geq1}5^rR_1=0.
\]
But direct substitution in the Weierstrass invariant gives
\begin{align*}
 j(A_{U_1})
 &=\frac{6912}{4+27(1+\epsilon)^2}
   =\frac{6912}{31+54\epsilon}\\
 &=\frac{6912}{31}-\frac{373248}{961}\epsilon,
\end{align*}
whose $\epsilon$ coefficient is nonzero. The map $R_1\to R_1[1/5]$ is injective, so this contradicts $j(A)=j(E)$.

The finite surjective cover $\Spec\mathbf Z_5\to S_1$ verifies the first hypothesis. Strict localization of the reduced trait has connected punctured spectrum. Nilpotent invariance of the \'etale site and the local-cohomology argument in Theorem~\ref{thm:positive} verify the depth hypothesis.
\end{proof}

On the characteristic-zero overlap, the $p$-divisible groups forget this infinitesimal change of $j$. Equality of those groups therefore provides no abelian-scheme gluing isomorphism. The next construction places the same discrepancy over a reduced base.

\section{A derivation subring with a dual-number completion}

We give the required case of the Ferrand--Raynaud--Olberding construction explicitly, including its noetherianity and normalization properties.

\begin{lemma}\label{lem:ring}
Let $V$ be a complete discrete valuation ring of mixed characteristic $(0,p)$, with uniformizer $p$, fraction field $K$, and residue field $k$. Let $C\subset K$ be a subfield for which $C\cap V\to k$ is surjective. Choose $J\in V$ transcendental over $C$, and let $F\subset K$ be algebraic over $C(J)$. Extend the derivation $D(J)=1$, $D(C)=0$ to $F$, and put
\[
 V_0=F\cap V,\qquad R=\{a\in V_0:D(a)\in V_0\}.
\]
Then $R$ is a one-dimensional noetherian local domain with fraction field $F$ and residue field $k$. Its normalization is the discrete valuation ring $V_0$, and $R\subset V_0$ induces identical fraction and residue fields. The normalization is not finite. The graph homomorphism
\begin{equation}\label{eq:graph}
 R\longrightarrow V[\epsilon]/(\epsilon^2),
 \qquad a\longmapsto a+\epsilon D(a)
\end{equation}
identifies the maximal-ideal completion of $R$ with $V[\epsilon]/(\epsilon^2)$. In addition,
\begin{equation}\label{eq:special}
 R/pR\xrightarrow{\sim}k[\epsilon]/(\epsilon^2),
 \qquad a\longmapsto\bar a+\epsilon\overline{D(a)}.
\end{equation}
\end{lemma}
\begin{proof}
Write $v$ for the valuation on $K$, normalized by $v(p)=1$, and set $v(0)=+\infty$. In characteristic zero every algebraic extension is separable, so the derivation extends uniquely: differentiating a minimal polynomial determines the derivative of an algebraic element, since its polynomial derivative is nonzero. The restriction of the valuation of $V$ to $F$ has value group $\mathbf Z$, as $p\in C\subset F$. Thus $V_0$ is a discrete valuation ring with uniformizer $p$ and residue field $k$.

Multiplication by a sufficiently high power of $p$ makes both an element of $F$ and its derivative integral. Therefore $R[1/p]=F$. An element $a\in R$ of valuation zero has inverse in $R$, because $D(a^{-1})=-a^{-2}D(a)$. Hence $R$ is local with maximal ideal
\[
 \mathfrak m=\{a\in R:v(a)>0\},
\]
and has residue field $k$.

Choose $c\in C\cap V$ with residue $\bar J$, and set $z=J-c$. Then $z\in\mathfrak m$ and $D(z)=1$. For $a\in\mathfrak m$, choose $b\in C\cap V$ with residue $\overline{D(a)}$. Both $a-bz$ and $D(a-bz)$ lie in $pV_0$, so $(a-bz)/p\in R$. Consequently
\begin{equation}\label{eq:mpowers}
 \mathfrak m=(p,z),\qquad
 \mathfrak m^2=p\mathfrak m,\qquad
 \mathfrak m^r=p^{r-1}\mathfrak m\quad(r\geq1).
\end{equation}
For the middle equality, the only additional generator to check is $z^2$; the element $z^2/p$ has positive valuation and derivative $2z/p\in V_0$, so belongs to $\mathfrak m$.

Equivalently,
\[
 \mathfrak m^r=\{a\in F:v(a)\geq r,\ v(Da)\geq r-1\}.
\]
Let $0\ne a\in\mathfrak m$ and write $s=v(a)\geq1$. If $b\in\mathfrak m^{2s}$, then $b/a\in V_0$, and
\[
 D(b/a)=\frac{aD(b)-bD(a)}{a^2}\in V_0:
\]
the two numerator terms have valuations at least $3s-1$ and $2s$ respectively. Thus $\mathfrak m^{2s}\subset aR$. Since $\mathfrak m$ is finitely generated, $R/\mathfrak m^{2s}$ has finite length, and so does $R/aR$. Every nonzero proper ideal contains such an $a$; lifting finitely many generators from $R/aR$ shows that it is finitely generated. Hence $R$ is noetherian. Every nonzero prime contains a power of $\mathfrak m$, so $\dim R=1$.

For $b\in V_0$, choose $N$ so large that
\[
 D(b^{p^N})=p^N b^{p^N-1}D(b)\in V_0.
\]
Then $b^{p^N}\in R$, proving that $b$ is integral over $R$. Conversely, an element of $F$ integral over $R$ is integral over $V_0$ and therefore belongs to $V_0$. This proves the normalization assertion. There are exactly two prime ideals in either ring, the generic and closed points, and the residue fields at both agree. Thus the normalization map is a universal homeomorphism; in particular $R$ is geometrically unibranch.

The ring $C\cap V$ is dense in $V$: recursively lift a residue, subtract it, divide the remainder by $p$, and repeat. Choose $c_N\in C\cap V$ with $v(J-c_N)\geq N$. The image of \eqref{eq:graph} is dense for the $(p,\epsilon)$-adic topology, since elements
\[
 a_0+b_0(J-c_N),\qquad a_0,b_0\in C\cap V,
\]
have derivative $b_0$ and approximate any prescribed pair of coefficients in $V+\epsilon V$. Formula~\eqref{eq:mpowers} shows that the inverse image of
\[
 (p,\epsilon)^r=p^rV+\epsilon p^{r-1}V
\]
is exactly $\mathfrak m^r$. Thus the graph induces isomorphisms on every quotient and hence on completions. Reduction modulo $p$ has kernel $pR$, because $a,Da\in pV_0$ exactly when $a/p\in R$. Surjectivity in \eqref{eq:special} follows from the residue representatives and from $z\mapsto\epsilon$.

Finally, a finite birational extension $R\subset V_0$ would have a nonzero conductor: a common denominator of finitely many $R$-module generators would give $0\ne d\in R$ with $dV_0\subset R$. This is impossible. Put $u_N=(J-c_N)/p^N\in V_0$. Then
\[
 D(du_N)=D(d)u_N+d/p^N.
\]
For $N>v(d)$ the first summand is integral and the second has negative valuation; therefore $du_N\notin R$. The normalization is not finite.
\end{proof}

In the notation of \citet[Theorem~4.1]{OlberdingBad}, one may take the immediate extension of discrete valuation rings $C\cap V\subset V_0$ and the rank-one lattice $V_0\,dJ$ in $\Omega_{F/C}=F\,dJ$. The graph completion in Lemma~\ref{lem:ring} is the rank-one instance of his Theorem~4.4. The explicit proof above is included because the values of $D(J)$ and the special-fibre map \eqref{eq:special} carry the elliptic contradiction.

\section{An ordinary canonical lift with a moving special fibre}

We now choose the ring data so that a split ordinary $p$-divisible group exists generically while the derivative of its elliptic $j$-invariant is nonzero.

Let $k=\overline{\mathbf F_5(t)}$ and $V=W(k)$, with fraction field $K$. The elliptic curve
\begin{equation}\label{eq:model}
 E_t:\quad y^2=x^3-\frac{3t}{t-1728}x-\frac{2t}{t-1728}
\end{equation}
over $k$ has $j$-invariant $t$. It is ordinary: in characteristic $5$, the Hasse coefficient of $y^2=f(x)$ is the coefficient of $x^4$ in $f(x)^2$. Here that coefficient is
\[
 -\frac{6t}{t-1728}=\frac{4t}{t-3}\ne0.
\]
Equivalently, the Cartier operator on its one-dimensional space of regular differentials is nonzero.

Serre--Tate theory gives its canonical ordinary lift $E^{\mathrm{can}}/V$, for which
\[
 E^{\mathrm{can}}[5^\infty]
 \simeq\mu_{5^\infty,V}\oplus(\mathbf Q_5/\mathbf Z_5)_V.
\]
For completeness, the formal lift is the compatible system corresponding to the split group over every $V/5^r$. In dimension one the line bundle of the divisor three times the zero section is compatible and relatively ample. Grothendieck's existence theorem algebraizes the system and that line bundle \citep[Tag~089A]{Stacks}; smoothness and the group laws follow by formal completion. This is the canonical elliptic lift used here. Set $J=j(E^{\mathrm{can}})$. Its residue is $t$, so $J$ is transcendental over $\mathbf Q$: if $J$ were algebraic over $\mathbf Q$, a primitive integral polynomial for $J$, reduced modulo $5$ after removing its common power of $5$, would give a nonzero polynomial over $\mathbf F_5$ satisfied by $t$.

The countable field $k$ admits an enumeration $a_1,a_2,\ldots$. Each residue class in $V$ is uncountable. Choose recursively $c_i\in V$ with residue $a_i$ which is transcendental over $\mathbf Q(J,c_1,\ldots,c_{i-1})$. That field and its algebraic closure in $K$ are countable, so this choice is possible at each step. The elements $J,c_1,c_2,\ldots$ are algebraically independent over $\mathbf Q$. Consequently
\[
 C=\mathbf Q(c_1,c_2,\ldots)
\]
has $C\cap V\to k$ surjective and leaves $J$ transcendental over $C$.

Replace $t$ by $J$ in \eqref{eq:model} to obtain $E_J$. Direct calculation gives
\begin{equation}\label{eq:invariants}
 \Delta(E_J)=\frac{1728^2J^2}{(J-1728)^3},
 \qquad j(E_J)=J.
\end{equation}
Both $J$ and $J-1728$ are units. Thus $E_J$ is an elliptic scheme over $V$. It is isomorphic over $V$ to $E^{\mathrm{can}}$: away from $j=0,1728$, elliptic curves with the same $j$ differ by a quadratic twist, and every unit of the complete strictly henselian ring $V$ is a square, by Hensel's lemma and $2\in V^\times$. Fix one resulting isomorphism of its $5$-divisible group with the split group.

Let $F$ be the relative algebraic closure of $C(J)$ in $K$. Every finite-level generic isomorphism just chosen is defined over $F$. Indeed, over characteristic zero, the isomorphism scheme between the two finite \'etale $5^r$-torsion group schemes is finite \'etale; each of its coordinates in $K$ is algebraic over $C(J)$. Compatibility of these finite-level maps is preserved on restriction to $F$.

Apply Lemma~\ref{lem:ring} with this field $F$ and the derivation $D(J)=1$, $D(C)=0$. Since $J\in R$ and $J-1728$ is a unit of $R$, the equation for $E_J$ is already an elliptic scheme over $R$, with discriminant a unit by \eqref{eq:invariants}. Set
\[
 A_U=(E_J)_F,\qquad
 M=\mu_{5^\infty,R}\oplus(\mathbf Q_5/\mathbf Z_5)_R,
\]
and use the inverse of the chosen generic splitting as $\varphi:M_F\simeq A_U[5^\infty]$.

\begin{proof}[Proof of Theorem~\ref{thm:reduced}]
All assertions about $R$, its dimension and its normalization follow from Lemma~\ref{lem:ring}. The finite-cover hypothesis holds with the identity because $R$ is geometrically unibranch. Here is also a direct verification of the depth condition for the chosen $F$. The discrete valuation ring $V_0=F\cap V$ is henselian: a simple root modulo its maximal ideal lifts to $V$ by Hensel's lemma, and the lift is algebraic over $F$, hence lies in $F$ by relative algebraic closedness. Its residue field is algebraically closed, so $V_0$ is strictly henselian. Since $\Spec V_0\to\Spec R$ is a universal homeomorphism, their \'etale sites agree \citep[Tag~04DY]{Stacks}. In particular $R$ is strictly henselian, and its punctured spectrum is the single generic point. The local-cohomology argument of Section~2 proves \eqref{eq:depth}.

Suppose the prescribed triple has a realization $A/R$. Its relative dimension is one. The generic isomorphism with $A_U$ gives $j(A)_F=J$. Because $R\hookrightarrow F$ is injective and $j$ is a regular invariant of an elliptic scheme, $j(A)=J$ in $R$. Under \eqref{eq:special}, this becomes
\[
 j(A\bmod5)=t+\epsilon.
\]
On the other hand, the prescribed group $M\bmod5$ over $k[\epsilon]/(\epsilon^2)$ is the constant split group. Let $A_0/k$ be the special fibre. The given group identification and its reduction identify this deformation datum with the constant deformation of $A_0[p^\infty]$. Serre--Tate over the square-zero extension $k\subset k[\epsilon]/(\epsilon^2)$ therefore gives an isomorphism between $A\bmod5$ and the constant elliptic deformation of $A_0$, respecting its reduction. Its $j$-invariant is $t$. This contradicts $t+\epsilon\ne t$, proving nonrealizability. The elliptic scheme $E_J/R$ already constructed proves the additional good-model assertion.
\end{proof}

\section{What the examples decide}

The positive characteristic example and the two mixed-characteristic obstructions isolate three different phenomena. If $p$ is locally nilpotent, a fixed reduced realization lifts with its prescribed group by Serre--Tate. If $p$ is invertible on the overlap, the torsion group can forget a first-order change of the abelian scheme. Finally, even a reduced one-dimensional geometrically unibranch base can have a nonreduced characteristic-$p$ fibre and completion that expose exactly this discrepancy.

These conclusions do not classify the essential image of $F_{S,U}$ for every nonnormal base. In particular, the hypotheses under which finite normalization suffices for a given gluing argument require a separate descent theorem with its precise compatibility assumptions. The obstruction exhibited here is concrete: a chosen generic $p$-divisible comparison need not be induced by an isomorphism of abelian schemes on the overlap of the open part and a formal model. A general criterion must retain that abelian-scheme comparison, rather than infer it from the group comparison alone.

\section*{Note on method and records}

The mathematics and text were produced by GPT-6 Astra Pro in ChatGPT Pro mode. Earlier targeted proof checks were recorded by Claude Opus~5.5 at effort max. A separate instance of the same model reconstructed the ring lemma and both main examples from statements and definitions alone, before reading the earlier proofs or any construction hint. Its published inputs were then checked at the exact statements, and its proofs were compared with the arguments here. Two earlier derivations had a short description of the infinitesimal strategy available and are recorded separately as targeted checks. The finite Weierstrass calculations were recomputed by independent exact polynomial arithmetic, and the original symbolic calculation was rerun. The program \emph{Elliptic invariant checks} and its expected output are available through the workbench. The published inputs were read at the stated locations. No human has checked the mathematics. The text is dedicated to the public domain under CC0~1.0. Sources and verification records belong to \url{\workbenchurl}; their permanent records address is \url{\recordurl}.

\begin{thebibliography}{99}
\bibitem[de Jong(1998)]{DeJongICM}
de Jong, A. J. (1998). Barsotti--Tate groups and crystals. In \emph{Proceedings of the International Congress of Mathematicians 1998, II} (pp.~259--265). Documenta Mathematica, Extra Volume. \url{https://doi.org/10.4171/DMS/1-2/25}
\bibitem[Ferrand and Raynaud(1970)]{FerrandRaynaud}
Ferrand, D., \& Raynaud, M. (1970). Fibres formelles d'un anneau local noeth\'erien. \emph{Annales scientifiques de l'\'Ecole Normale Sup\'erieure, 4e s\'erie, 3}(3), 295--311. \url{https://www.numdam.org/item/ASENS_1970_4_3_3_295_0/}
\bibitem[Grothendieck(n.d.)]{GrothendieckNotes}
Grothendieck, A. (n.d.). \emph{Bonne r\'eduction des vari\'et\'es ab\'eliennes via bonne r\'eduction de Barsotti--Tate}. Archives Grothendieck, Universit\'e de Montpellier, cote~100. \url{https://grothendieck.umontpellier.fr/100.pdf}
\bibitem[Grothendieck(2005)]{SGA2}
Grothendieck, A. (2005). \emph{Cohomologie locale des faisceaux coh\'erents et th\'eor\`emes de Lefschetz locaux et globaux (SGA~2)}. Documents Math\'ematiques, Vol.~4. Soci\'et\'e Math\'ematique de France. Expos\'e~XIV by M. Raynaud, Remark~1.6(b), p.~164. \url{https://www.cmls.polytechnique.fr/perso/laszlo/sga2/sga2-smf.pdf}
\bibitem[Katz(1981)]{Katz}
Katz, N. M. (1981). Serre--Tate local moduli. In \emph{Surfaces alg\'ebriques} (Lecture Notes in Mathematics, Vol.~868, pp.~138--202). Springer. \url{https://web.math.princeton.edu/~nmk/old/serretatelocmod.pdf}
\bibitem[Krishnamoorthy and P\'al(2021)]{KrishnamoorthyPal}
Krishnamoorthy, R., \& P\'al, A. (2021). Rank~2 local systems and abelian varieties. \emph{Selecta Mathematica, 27}, Article~51. \url{https://doi.org/10.1007/s00029-021-00669-8}
\bibitem[Liu and Tong(2016)]{LiuTong}
Liu, Q., \& Tong, J. (2016). N\'eron models of algebraic curves. \emph{Transactions of the American Mathematical Society, 368}(10), 7019--7043. \url{https://doi.org/10.1090/tran/6642} Author manuscript: \url{https://www.math.u-bordeaux.fr/~qliu/articles/Neron_models_curves.pdf}
\bibitem[Olberding(2012)]{OlberdingBad}
Olberding, B. (2012). \emph{One-dimensional bad Noetherian domains}. arXiv:1209.2991v1. \url{https://arxiv.org/abs/1209.2991}
\bibitem[The Stacks Project Authors(2026)]{Stacks}
The Stacks Project Authors. (2026). \emph{The Stacks Project}. Tag~04DY: Topological invariance of the small \'etale site; Tag~089A: Grothendieck's algebraization theorem. \url{https://stacks.math.columbia.edu/tag/04DY}
\end{thebibliography}
\end{document}
